\input{../tex-inputs/latex-korrekturansicht-vorspann}

\section[Arthur Schnitzler an Theodor Herzl, 28. 9. 1893]{L03906 Arthur Schnitzler an Theodor Herzl, 28. 9. 1893}
\nopagebreak\mylabel{L03906v}
\rehead{ }\normalsize\beginnumbering\briefempfaengerindex{Herzl, Theodor@\textsc{Herzl, Theodor}!zzzSchnitzler, Arthur@\emph{von Arthur Schnitzler}!1893-09-281@{28. 9. 1893}|(be}
\toendnotes[C]{\smallbreak\pagebreak[2]}
\correspDesc{Versand  durch Arthur Schnitzler am 28. 9. 1893 in Wien
\newline{}Erhalt  durch Theodor Herzl im Zeitraum [28. 9. 1893
                  – 1. 10. 1893?] in Wien}\toendnotes[C]{\smallbreak}
\Standort{Jerusalem, Central Zionist Archives, H1:1924-11.}
\physDesc{Brief, 1 Blatt, 4 Seiten, 1251 Zeichen (Briefpapier mit Trauerrand)
\newline{}Handschrift: schwarze Tinte, deutsche Kurrent
\newline{}Ordnung: mit Bleistift von unbekannter Hand innerhalb das Konvoluts paginiert:
                                    »35«–»36« }
\buchAbdrucke{\weitereDrucke{\emph{2. [10.] 93.} In: Arthur Schnitzler: \emph{Briefe 1875–1912}. Herausgegeben von Therese Nickl und Heinrich Schnitzler. Frankfurt am Main: \emph{S. Fischer} 1981, S. 213–214.} }\toendnotes[C]{\smallbreak}
\pstart
           \raggedleft{}{\pb}\uline{28. 9. 93}\pend
           
\pstart{}Verehrter Freund!\pend\vspace{0.5em}
\pstart
           mit ½ 11 war ich geſtern Abend im \textcolor{pink}{\textsc{Meissl u Schadn}}\oindex{Wien@\textbf{Wien}!I., Innere Stadt@\textbf{I., Innere Stadt}!Meissl {\kaufmannsund} Schadn@\textbf{Meissl {\kaufmannsund} Schadn}, \emph{Hotel}|pw}{}\ledrightnote{\textcolor{pink}{Meissl {\kaufmannsund} Schadn}}; man hatte mir, \textsc{resp.} meiner \textcolor{blue}{Mama}\pwindex{Schnitzler, Louise 8.\,7.\,1840 Kőszeg – 9.\,9.\,1911 Wien@\textsc{Schnitzler, Louise} (8.\,7.\,1840 Kőszeg – 9.\,9.\,1911 Wien)|pwv}{}\ledrightnote{{$\rightarrow$}\emph{\textcolor{blue}{Louise Schnitzler}}} nicht telephonirt, \uline{wer}
               mich hinbeſtellt hatte; und ſo glaubte ich, da ich \label{K_L03906-1v}\edtext{beka{\geminationn}tlich ein beſchäftigter
               praktiſcher Arzt}{\lemma{\textnormal{\emph{bekanntlich … Arzt}}}\Cendnote{\textnormal{\textcolor{blue}{Schnitzler} hatte wenige Privatpatienten, so
                  dass hier ein ironischer Unterton vorliegt.}}}\label{K_L03906-1} bin, daſs es ſich um einen
               Patienten handle. Es wurde nun in allen Stöcken durch Portier, {\pb}Zi{\geminationm}erkellner u. Stubenmädchen
               herumgefragt, wer denn krank ſei; – es meldete ſich ſchließlich die Nu{\geminationm}er 19, die aber von mir nichts wiſſen wollte. Ich ging
               ſchwer gekränktens davon, dachte an das Bubenstück irgend eines Neiders. Ko{\geminationm} ins \textcolor{pink}{Grienſteidl}\oindex{Wien@\textbf{Wien}!I., Innere Stadt@\textbf{I., Innere Stadt}!Café Griensteidl@\textbf{Café Griensteidl}, \emph{Kaffeehaus}|pw}{}\ledrightnote{\textcolor{pink}{Café Griensteidl}}. Wie ich, ſo gegen vier, im Weggehn
               bin, erzählt mir der \textcolor{blue}{Heinrich}\pwindex{Heinrich @\textsc{Heinrich}, \emph{Kellner}|pw}{}\ledrightnote{\textcolor{blue}{Heinrich}}, Sie wären da
               geweſen, u er {\pb}hätte Ihre Poſt durch einen Burſchen an den
                  \textcolor{pink}{Burgring}\oindex{Wien@\textbf{Wien}!I., Innere Stadt@\textbf{I., Innere Stadt}!Wohnung und Ordination Johann Schnitzler Burgring 1@\textbf{Wohnung und Ordination Johann Schnitzler Burgring 1}, \emph{Ordination}|pw}{}\ledrightnote{\textcolor{pink}{Wohnung und Ordination Johann Schnitzler Burgring 1}} telephoniren laſſen. Auf dieſem
               Wege iſt nun nichts geringeres als Ihr Name in Verluſt gerathen. –\pend
           
\pstart
           Dieser einfachen Geſchichte hab ich nur mehr hinzuzuſetzen, daſs ich untröſtlich bin,
               aber dennoch hoffe, Sie vor Ihrer Abreiſe noch zu ſehn? – So gegen {\pb}eilf bin ich wohl alle Abende im \textsc{\textcolor{pink}{Griensteidl}\oindex{Wien@\textbf{Wien}!I., Innere Stadt@\textbf{I., Innere Stadt}!Café Griensteidl@\textbf{Café Griensteidl}, \emph{Kaffeehaus}|pw}{}\ledrightnote{\textcolor{pink}{Café Griensteidl}}} was ich Ihnen für den Fall eines Theaterbeſuchs mittheile, Samſtag
               geh’ ich \textcolor{violet}{zur \textcolor{green}{Sündfluth}\pwindex{Brociner, Marco 20.\,10.\,1852 Iași – 12.\,4.\,1942 Wien@\textsc{Brociner, Marco} (20.\,10.\,1852 Iași – 12.\,4.\,1942 Wien), \emph{Schriftsteller, Journalist, Kritiker}!Sündfluth. Schauspiel in vier Akten@\strich\emph{Die Sündfluth. Schauspiel in vier Akten}|pw}\pwindex{Loewe, Konrad 6.\,2.\,1856 Prostějov – 11.\,2.\,1912 Wien@\textsc{Loewe, Konrad} (6.\,2.\,1856 Prostějov – 11.\,2.\,1912 Wien), \emph{Schriftsteller, Schauspieler}!Sündfluth. Schauspiel in vier Akten@\strich\emph{Die Sündfluth. Schauspiel in vier Akten}|pw}{}\ledrightnote{\textcolor{green}{Die Sündfluth. Schauspiel in vier Akten}}}\eventindex{Volkstheater@\textbf{Volkstheater}!Premiere von Die Sündfluth, 7.10.1893@Premiere von Die Sündfluth, 7.10.1893|pw}{}\ledrightnote{\textcolor{violet}{Premiere von Die Sündfluth, 7.10.1893}}.\pend
           
\pstart
           – Also auf Wiederſehen! Empfehlen Sie mich zutiefſt Ihrer w. Frau \textcolor{blue}{Gemahlin}\pwindex{Herzl, Julie 1.\,2.\,1868 Budapest – 10.\,11.\,1907 Wien@\textsc{Herzl, Julie} (1.\,2.\,1868 Budapest – 10.\,11.\,1907 Wien)|pwv}{}\ledrightnote{{$\rightarrow$}\emph{\textcolor{blue}{Julie Herzl}}} und allen übrigen Mitgliedern Ihrer
                  liebenswürdig\textcolor{gray}{en} Familie.\pend
           
\pstart
           Herzlichen Gruſs{\\[\baselineskip]}Ihr getreuer \spacefill\mbox{ArthurSchnitzler}\pend
           \leftskip=0em{}\selectlanguage{ngerman}\endnumbering\briefempfaengerindex{Herzl, Theodor@\textsc{Herzl, Theodor}!zzzSchnitzler, Arthur@\emph{von Arthur Schnitzler}!1893-09-281@{28. 9. 1893}|)be}\mylabel{L03906h}  \newcommand{\dateiname}{L03906}\newcommand{\titel}{Arthur Schnitzler an Theodor Herzl, 28. 9. 1893}\newcommand{\editorInnen}{Herausgegeben von Jahnke, SelmaMüller, Martin Anton}\input{../tex-inputs/latex-korrekturansicht-abspann}
